\documentclass[aspectratio=169,11pt]{beamer}
\usepackage{fontspec}
\usepackage{listings}
\usepackage{xcolor}
\usepackage{booktabs}
\usepackage{array}

% ================= Palette (identique au livre Tome 1) =================
\definecolor{PyBleu}{HTML}{1A4A7A}
\definecolor{PyBleuMoyen}{HTML}{2D7AB3}
\definecolor{PyBleuClair}{HTML}{EAF1F8}
\definecolor{PyJaune}{HTML}{FFD43B}
\definecolor{PyNuit}{HTML}{142840}
\definecolor{PyTexte}{HTML}{1F2A37}
\definecolor{PyGris}{HTML}{5A6B7E}
\definecolor{PyRouge}{HTML}{B3442D}
\definecolor{PyVert}{HTML}{2E7D5B}
\definecolor{CodeTexte}{HTML}{E8EEF6}

% ================= Code Python =================
\lstdefinestyle{py}{
  language=Python,
  keywordstyle=\color{PyJaune}\bfseries,
  stringstyle=\color{PyJaune},
  commentstyle=\color{PyGris}\itshape,
  showstringspaces=false,
  breaklines=true,
  backgroundcolor=\color{PyNuit},
  basicstyle=\ttfamily\footnotesize\color{CodeTexte},
  columns=fullflexible,
  keepspaces=true,
  xleftmargin=6pt,
  framexleftmargin=6pt,
}
\lstdefinestyle{pysmall}{
  language=Python,
  keywordstyle=\color{PyJaune}\bfseries,
  stringstyle=\color{PyJaune},
  commentstyle=\color{PyGris}\itshape,
  showstringspaces=false,
  breaklines=true,
  backgroundcolor=\color{PyNuit},
  basicstyle=\ttfamily\scriptsize\color{CodeTexte},
  columns=fullflexible,
  keepspaces=true,
  xleftmargin=5pt,
  framexleftmargin=5pt,
}

% ================= Titres de diapositives =================
\setbeamercolor{structure}{fg=PyBleu}
\setbeamerfont{frametitle}{series=\bfseries, size=\Large}
\setbeamercolor{frametitle}{bg=PyBleu, fg=white}
\setbeamertemplate{frametitle}{%
  \nointerlineskip
  \begin{beamercolorbox}[wd=\textwidth, ht=2.1ex, dp=0.9ex,
                          leftskip=0.35cm, rightskip=0.35cm]{frametitle}%
    \usebeamerfont{frametitle}\insertframetitle%
  \end{beamercolorbox}%
  \vspace*{-0.4pt}%
  {\color{PyJaune}\rule{\textwidth}{1.6pt}}%
}

% ================= Pied de page =================
\setbeamerfont{footline}{size=\tiny}
\setbeamercolor{footline}{bg=white, fg=PyGris}
\newcommand{\PiedDePage}{Python en 24 heures \textbar\ Séance 05 : Listes, tuples et compréhensions}
\setbeamertemplate{footline}{%
  \leavevmode%
  \hbox{%
    \begin{beamercolorbox}[wd=\paperwidth, ht=2.2ex, dp=1.1ex,
                            leftskip=0.35cm, rightskip=0.35cm]{footline}%
      \usebeamerfont{footline}\PiedDePage
      \hfill\insertframenumber\,/\,\inserttotalframenumber%
    \end{beamercolorbox}}%
}

\setbeamertemplate{navigation symbols}{}
\setbeamertemplate{itemize item}{\color{PyBleuMoyen}$\blacktriangleright$}
\setbeamertemplate{itemize subitem}{\color{PyBleuMoyen}\textbullet}

% ================= Blocs pedagogiques =================
\setbeamercolor{block title}{bg=PyBleu, fg=white}
\setbeamercolor{block body}{bg=PyBleuClair, fg=PyTexte}
\setbeamercolor{block title alerted}{bg=PyRouge, fg=white}
\setbeamercolor{block body alerted}{bg=PyRouge!9, fg=PyTexte}
\setbeamercolor{block title example}{bg=PyJaune, fg=PyNuit}
\setbeamercolor{block body example}{bg=PyJaune!14, fg=PyTexte}

% ================= Metadonnees =================
\hypersetup{
  pdftitle={Python en 24 heures : Séance 05 : Listes, tuples et compréhensions},
  pdfauthor={HOUNKPE Mawulikplimi Roland et DJABON Ounimborbitibou},
  pdfsubject={De l'algorithmique à la programmation orientée objet},
  pdfkeywords={Python, algorithmique, POO, formation, exercices corrigés}
}

\begin{document}

{%
\setbeamercolor{background canvas}{bg=PyNuit}
\begin{frame}[plain]
\vspace{10mm}
{\color{PyJaune}\small\bfseries FORMATION « PYTHON EN 24 HEURES » \quad\textbar\quad SÉANCE 05\par}
\vspace{4mm}
{\color{white}\Huge\bfseries Listes, tuples et compréhensions\par}
\vspace{4mm}
{\color{PyBleuClair}\large Chapitres 9 et 10 du Tome 1\par}
\vspace{7mm}
{\color{PyJaune}\rule{34mm}{2.2pt}\par}
\vspace{7mm}
{\color{white}\normalsize 1 h de cours (présentation Beamer) \quad$\bullet$\quad 1 h de pratique (notebook Jupyter)\par}
\vspace{2.5mm}
{\color{PyBleuClair}\normalsize HOUNKPE Mawulikplimi Roland et DJABON Ounimborbitibou\par}
\end{frame}
}
\begin{frame}{Objectifs de la séance}
\begin{itemize}
  \item Créer, lire et modifier des listes avec index, slices et méthodes
  \item Parcourir une liste avec for, enumerate et le test d'appartenance in
  \item Distinguer les tuples, figés une fois pour toutes, et leurs usages
  \item Écrire des compréhensions de listes simples, filtrées et imbriquées
  \item Éviter le piège de la copie de liste partagée
\end{itemize}
\end{frame}
\begin{frame}{Rappel de la séance 04}
Les listes sont le premier conteneur du cours : une seule variable pour ranger toute une série, parcourue par les boucles que vous maîtrisez déjà.

\begin{itemize}
  \item if, elif, else : les décisions en Python et l'indentation structurante
  \item and, or, not : les connecteurs logiques
  \item while et for : les deux boucles, break et continue
  \item range(a, b) compte de a à b moins 1 : la borne de fin est exclue
  \item Une comparaison produit un bool : c'est le carburant des filtres
\end{itemize}
\end{frame}
\begin{frame}[fragile]{Créer et lire une liste}

\begin{lstlisting}[style=py]
notes = [12, 15, 9, 18]
print(notes[0])      # 12  premier élément
print(notes[-1])     # 18  dernier élément
print(notes[1:3])    # [15, 9]  slice : borne de fin exclue
print(len(notes))    # 4
\end{lstlisting}
\vspace{0.5pt}

{\small\color{PyGris}\itshape Les index démarrent à 0. Les index négatifs comptent depuis la fin : notes[-1] est le dernier.}

\end{frame}

\begin{frame}[fragile]{Modifier une liste : les méthodes clés}
\begin{columns}[t]
\begin{column}{0.42\textwidth}
\begin{itemize}
  \setlength{\itemsep}{1.5pt}\setlength{\parskip}{0pt}
  \item `append(x)` : ajoute à la fin
  \item `insert(i, x)` : insère à la position i
  \item `remove(x)` : retire la première occurrence de x
  \item `pop(i)` : retire et \textbf{retourne} l'élément à l'index i
  \item `sort()` et `reverse()` : rangent la liste \textbf{sur place}
  \item `in` : teste l'appartenance, renvoie un bool
\end{itemize}
\end{column}
\hspace{0.02\textwidth}
\begin{column}{0.55\textwidth}
\begin{lstlisting}[style=pysmall]
panier = ["riz", "huile"]
panier.append("sucre")
panier.insert(0, "pain")
panier.remove("huile")
dernier = panier.pop()
print(panier, dernier)
print("riz" in panier)
\end{lstlisting}
\end{column}
\end{columns}
\end{frame}

\begin{frame}[fragile]{Parcourir : for, enumerate}

\begin{lstlisting}[style=py]
notes = [12, 15, 9]
for n in notes:
    print(n)

for i, n in enumerate(notes):
    print(f"Élément {i} : {n}")
\end{lstlisting}
\vspace{0.5pt}

{\small\color{PyGris}\itshape Préférez le parcours direct for n in notes à la manipulation d'index. enumerate fournit la position quand elle est vraiment utile.}

\end{frame}

\begin{frame}[fragile]{Le tuple : une liste figée}
\begin{columns}[t]
\begin{column}{0.42\textwidth}
\begin{itemize}
  \setlength{\itemsep}{1.5pt}\setlength{\parskip}{0pt}
  \item Parenthèses au lieu de crochets : (3, 5)
  \item \textbf{Immuables} : impossible de modifier un élément
  \item Ideal pour un résultat composé : (somme, moyenne)
  \item Unpacking : `a, b = point` déplie le tuple en deux variables
  \item Retourner plusieurs valeurs d'une fonction : c'est un tuple
\end{itemize}
\end{column}
\hspace{0.02\textwidth}
\begin{column}{0.55\textwidth}
\begin{lstlisting}[style=pysmall]
point = (3, 5)
x, y = point          # unpacking
print(x + y)

def stats(notes):
    return (min(notes), max(notes))
mini, maxi = stats([7, 9, 14])
print(mini, maxi)
\end{lstlisting}
\end{column}
\end{columns}
\end{frame}

\begin{frame}{La compréhension de liste : l'idée}
\begin{itemize}
  \item Un motif ultra courant : créer une liste en transformant une autre suite
  \item Version classique : boucle for + append, trois lignes
  \item La compréhension tient en \textbf{une seule ligne lisible}
  \item `[expression for element in suite]`
  \item Python exécute exactement la même chose, mais le code dit l'essentiel en premier
\end{itemize}
\end{frame}

\begin{frame}[fragile]{Boucle classique contre compréhension}

\begin{lstlisting}[style=py]
# Version classique
cubes = []
for k in range(1, 6):
    cubes.append(k ** 3)

# Même résultat en compréhension
cubes = [k ** 3 for k in range(1, 6)]
print(cubes)   # [1, 8, 27, 64, 125]
\end{lstlisting}
\vspace{0.5pt}

{\small\color{PyGris}\itshape On lit de gauche à droite : « le cube de k, pour k allant de 1 à 5 ». L'ordre d'écriture suit l'ordre de lecture du résultat.}

\end{frame}

\begin{frame}[fragile]{Compréhension avec condition}

\begin{lstlisting}[style=py]
notes = [12, 15, 9, 18, 8, 14]

admis = [n for n in notes if n >= 10]
print(admis)          # [12, 15, 18, 14]

recales = [n for n in notes if n < 10]
print(recales)        # [9, 8]
\end{lstlisting}
\vspace{0.5pt}

{\small\color{PyGris}\itshape Le `if` en fin de compréhension filtre : seuls les éléments qui passent le test entrent dans la liste.}

\end{frame}

\begin{frame}[fragile]{Transformer et filtrer en même temps}

\begin{lstlisting}[style=py]
prix_ht = [100, 250, 40, 999]

prix_ttc = [round(p * 1.18, 2) for p in prix_ht]
print(prix_ttc)

gros_montants = [p * 1.18 for p in prix_ht if p >= 100]
print(gros_montants)
\end{lstlisting}
\vspace{0.5pt}

{\small\color{PyGris}\itshape L'expression de gauche peut calculer (arrondi, conversion) pendant que le if de droite filtre : les deux rôles sont séparés.}

\end{frame}

\begin{frame}[fragile]{Compréhensions imbriquées}

\begin{lstlisting}[style=py]
matrice = [[1, 2, 3],
           [4, 5, 6]]
ligne_plate = [v for ligne in matrice for v in ligne]
print(ligne_plate)    # [1, 2, 3, 4, 5, 6]

carrees = [(x, y) for x in range(3) for y in range(3)]
print(len(carrees))   # 9
\end{lstlisting}
\vspace{0.5pt}

{\small\color{PyGris}\itshape Deux for se lisent comme deux boucles imbriquées : le for de droite est la boucle interne. L'ordre des clauses est celui des boucles.}

\end{frame}

\begin{frame}[fragile]{Le piège de la copie}
\begin{columns}[t]
\begin{column}{0.42\textwidth}
\begin{itemize}
  \setlength{\itemsep}{1.5pt}\setlength{\parskip}{0pt}
  \item `b = a` ne copie \textbf{pas} la liste : les deux noms pointent vers le même objet
  \item Modifier b modifie a : surprise garantie
  \item Copie superficielle : `b = a.copy()` ou `b = list(a)` ou `b = a[:]`
  \item Vérification : `a is b` dit si les deux noms désignent le même objet
\end{itemize}
\end{column}
\hspace{0.02\textwidth}
\begin{column}{0.55\textwidth}
\begin{lstlisting}[style=pysmall]
a = [1, 2, 3]
b = a               # même objet !
b.append(4)
print(a)            # [1, 2, 3, 4]

c = a.copy()        # vraie copie
c.append(99)
print(a, c)
print(a is b, a is c)
\end{lstlisting}
\end{column}
\end{columns}
\end{frame}

\begin{frame}{Attention}
\begin{alertblock}{IndexError et slices}
`notes[10]` sur une liste de 4 éléments lève `IndexError` : l'index doit être strictement inférieur à la longueur. Étrangement, le slice `notes[10:20]` ne lève aucune erreur et retourne une liste vide. Rappel : `range(len(notes))` donne les index valides, mais le parcours direct `for n in notes` reste plus sûr.
\end{alertblock}
\end{frame}

\begin{frame}{À retenir}
\begin{block}{L'essentiel de la séance 05}
Une liste se parcourt directement et se modifie sur place avec append, insert, remove et pop. Le tuple fige des valeurs liées et sert à retourner plusieurs résultats. La compréhension de liste résume transformation et filtrage en une ligne : expression, for, condition. Et souvenez vous : affecter une liste à un second nom ne la copie pas.
\end{block}
\end{frame}

\begin{frame}{Travaux pratiques (1 h) : Listes, tuples et compréhensions}
Du panier de courses au nettoyage de données imbriquées : cinq exercices pour faire de la compréhension de liste un réflexe. Pour les exercices 2 à 5, une solution en une ligne de compréhension est attendue quand c'est possible.

\begin{itemize}
  \setlength{\itemsep}{1.5pt}\setlength{\parskip}{0pt}
  \item \textbf{Exercice 1 (Basique)} : Panier de courses
  \item \textbf{Exercice 2 (Basique)} : Carrés et cubes en compréhension
  \item \textbf{Exercice 3 (Intermédiaire)} : Filtrer les notes admises
  \item \textbf{Exercice 4 (Intermédiaire)} : Transposée d'une matrice
  \item \textbf{Exercice 5 (Défi)} : Nettoyage de données imbriquées
\end{itemize}
\vspace{4pt}

{\small\color{PyGris}Ouvrez le notebook « Seance\_05\_listes\_tuples\_comprehensions.ipynb » pour commencer.}
\end{frame}
\begin{frame}{Pour la prochaine séance}
\begin{exampleblock}{Conseil de travail}
Séance 06 : dictionnaires et ensembles pour organiser les données par clé, puis les fonctions pour organiser le code. Relisez le piège de la copie : il réapparaît avec les dictionnaires, exactement de la même façon.
\end{exampleblock}
\end{frame}

\end{document}
